% Options for packages loaded elsewhere
\PassOptionsToPackage{unicode}{hyperref}
\PassOptionsToPackage{hyphens}{url}
\PassOptionsToPackage{dvipsnames,svgnames,x11names}{xcolor}
\documentclass[
  turkish,
  10pt,
  a4paper,
]{article}
\usepackage{xcolor}
\usepackage[margin=20mm]{geometry}
\usepackage{amsmath,amssymb}
\setcounter{secnumdepth}{5}
\usepackage{iftex}
\ifPDFTeX
  \usepackage[T1]{fontenc}
  \usepackage[utf8]{inputenc}
  \usepackage{textcomp} % provide euro and other symbols
\else % if luatex or xetex
  \usepackage{unicode-math} % this also loads fontspec
  \defaultfontfeatures{Scale=MatchLowercase}
  \defaultfontfeatures[\rmfamily]{Ligatures=TeX,Scale=1}
\fi
\usepackage{lmodern}
\ifPDFTeX\else
  % xetex/luatex font selection
  \setmainfont[]{DejaVu Sans}
\fi
% Use upquote if available, for straight quotes in verbatim environments
\IfFileExists{upquote.sty}{\usepackage{upquote}}{}
\IfFileExists{microtype.sty}{% use microtype if available
  \usepackage[]{microtype}
  \UseMicrotypeSet[protrusion]{basicmath} % disable protrusion for tt fonts
}{}
\makeatletter
\@ifundefined{KOMAClassName}{% if non-KOMA class
  \IfFileExists{parskip.sty}{%
    \usepackage{parskip}
  }{% else
    \setlength{\parindent}{0pt}
    \setlength{\parskip}{6pt plus 2pt minus 1pt}}
}{% if KOMA class
  \KOMAoptions{parskip=half}}
\makeatother
\usepackage{color}
\usepackage{fancyvrb}
\newcommand{\VerbBar}{|}
\newcommand{\VERB}{\Verb[commandchars=\\\{\}]}
\DefineVerbatimEnvironment{Highlighting}{Verbatim}{commandchars=\\\{\}}
% Add ',fontsize=\small' for more characters per line
\newenvironment{Shaded}{}{}
\newcommand{\AlertTok}[1]{\textcolor[rgb]{1.00,0.00,0.00}{\textbf{#1}}}
\newcommand{\AnnotationTok}[1]{\textcolor[rgb]{0.38,0.63,0.69}{\textbf{\textit{#1}}}}
\newcommand{\AttributeTok}[1]{\textcolor[rgb]{0.49,0.56,0.16}{#1}}
\newcommand{\BaseNTok}[1]{\textcolor[rgb]{0.25,0.63,0.44}{#1}}
\newcommand{\BuiltInTok}[1]{\textcolor[rgb]{0.00,0.50,0.00}{#1}}
\newcommand{\CharTok}[1]{\textcolor[rgb]{0.25,0.44,0.63}{#1}}
\newcommand{\CommentTok}[1]{\textcolor[rgb]{0.38,0.63,0.69}{\textit{#1}}}
\newcommand{\CommentVarTok}[1]{\textcolor[rgb]{0.38,0.63,0.69}{\textbf{\textit{#1}}}}
\newcommand{\ConstantTok}[1]{\textcolor[rgb]{0.53,0.00,0.00}{#1}}
\newcommand{\ControlFlowTok}[1]{\textcolor[rgb]{0.00,0.44,0.13}{\textbf{#1}}}
\newcommand{\DataTypeTok}[1]{\textcolor[rgb]{0.56,0.13,0.00}{#1}}
\newcommand{\DecValTok}[1]{\textcolor[rgb]{0.25,0.63,0.44}{#1}}
\newcommand{\DocumentationTok}[1]{\textcolor[rgb]{0.73,0.13,0.13}{\textit{#1}}}
\newcommand{\ErrorTok}[1]{\textcolor[rgb]{1.00,0.00,0.00}{\textbf{#1}}}
\newcommand{\ExtensionTok}[1]{#1}
\newcommand{\FloatTok}[1]{\textcolor[rgb]{0.25,0.63,0.44}{#1}}
\newcommand{\FunctionTok}[1]{\textcolor[rgb]{0.02,0.16,0.49}{#1}}
\newcommand{\ImportTok}[1]{\textcolor[rgb]{0.00,0.50,0.00}{\textbf{#1}}}
\newcommand{\InformationTok}[1]{\textcolor[rgb]{0.38,0.63,0.69}{\textbf{\textit{#1}}}}
\newcommand{\KeywordTok}[1]{\textcolor[rgb]{0.00,0.44,0.13}{\textbf{#1}}}
\newcommand{\NormalTok}[1]{#1}
\newcommand{\OperatorTok}[1]{\textcolor[rgb]{0.40,0.40,0.40}{#1}}
\newcommand{\OtherTok}[1]{\textcolor[rgb]{0.00,0.44,0.13}{#1}}
\newcommand{\PreprocessorTok}[1]{\textcolor[rgb]{0.74,0.48,0.00}{#1}}
\newcommand{\RegionMarkerTok}[1]{#1}
\newcommand{\SpecialCharTok}[1]{\textcolor[rgb]{0.25,0.44,0.63}{#1}}
\newcommand{\SpecialStringTok}[1]{\textcolor[rgb]{0.73,0.40,0.53}{#1}}
\newcommand{\StringTok}[1]{\textcolor[rgb]{0.25,0.44,0.63}{#1}}
\newcommand{\VariableTok}[1]{\textcolor[rgb]{0.10,0.09,0.49}{#1}}
\newcommand{\VerbatimStringTok}[1]{\textcolor[rgb]{0.25,0.44,0.63}{#1}}
\newcommand{\WarningTok}[1]{\textcolor[rgb]{0.38,0.63,0.69}{\textbf{\textit{#1}}}}
\usepackage{longtable,booktabs,array}
\usepackage{calc} % for calculating minipage widths
% Correct order of tables after \paragraph or \subparagraph
\usepackage{etoolbox}
\makeatletter
\patchcmd\longtable{\par}{\if@noskipsec\mbox{}\fi\par}{}{}
\makeatother
% Allow footnotes in longtable head/foot
\IfFileExists{footnotehyper.sty}{\usepackage{footnotehyper}}{\usepackage{footnote}}
\makesavenoteenv{longtable}
\usepackage{graphicx}
\makeatletter
\newsavebox\pandoc@box
\newcommand*\pandocbounded[1]{% scales image to fit in text height/width
  \sbox\pandoc@box{#1}%
  \Gscale@div\@tempa{\textheight}{\dimexpr\ht\pandoc@box+\dp\pandoc@box\relax}%
  \Gscale@div\@tempb{\linewidth}{\wd\pandoc@box}%
  \ifdim\@tempb\p@<\@tempa\p@\let\@tempa\@tempb\fi% select the smaller of both
  \ifdim\@tempa\p@<\p@\scalebox{\@tempa}{\usebox\pandoc@box}%
  \else\usebox{\pandoc@box}%
  \fi%
}
% Set default figure placement to htbp
\def\fps@figure{htbp}
\makeatother
\ifLuaTeX
\usepackage[bidi=basic,shorthands=off,]{babel}
\else
\usepackage[bidi=default,shorthands=off,]{babel}
\fi
\ifPDFTeX
\else
\babelfont{rm}[]{DejaVu Sans}
\fi
\ifLuaTeX
  \usepackage{selnolig} % disable illegal ligatures
\fi
\setlength{\emergencystretch}{3em} % prevent overfull lines
\providecommand{\tightlist}{%
  \setlength{\itemsep}{0pt}\setlength{\parskip}{0pt}}
\usepackage{bookmark}
\IfFileExists{xurl.sty}{\usepackage{xurl}}{} % add URL line breaks if available
\urlstyle{same}
\hypersetup{
  pdftitle={Amarcode Staj Raporu - Ekler},
  pdflang={tr-TR},
  colorlinks=true,
  linkcolor={blue},
  filecolor={Maroon},
  citecolor={Blue},
  urlcolor={blue},
  pdfcreator={LaTeX via pandoc}}

\title{Amarcode Staj Raporu - Ekler}
\author{}
\date{}

\begin{document}
\maketitle

\renewcommand*\contentsname{İçindekiler}
{
\hypersetup{linkcolor=}
\setcounter{tocdepth}{3}
\tableofcontents
}
\section{Amarcode Staj Raporu ---
Ekler}\label{amarcode-staj-raporu--ekler}

Bu bölüm, 18 günlük staj raporunda açıklanan çalışmaları destekleyen
teknik materyalleri bir araya getirmektedir. Eklerde yer alan mimari
açıklamalar, tablolar ve görseller Amarcode kaynak kodu, Git geçmişi,
proje belgeleri ve uygulamanın son durumuna ait ekran görüntüleri temel
alınarak hazırlanmıştır.

\subsection{Ek A --- Proje Bileşenleri ve Depo
Yapısı}\label{ek-a--proje-bileux15fenleri-ve-depo-yapux131sux131}

Amarcode tek bir masaüstü arayüzünden oluşmamaktadır. Proje; kullanıcı
arayüzü, yerel masaüstü katmanı, bağımsız arka plan servisi, ortak
iletişim protokolü, ACP adaptörü ve kayıt hizmetinden meydana gelen bir
monorepo yapısına sahiptir.

\begin{longtable}[]{@{}lll@{}}
\toprule\noalign{}
Bileşen & Konum & Temel sorumluluk \\
\midrule\noalign{}
\endhead
\bottomrule\noalign{}
\endlastfoot
React uygulaması &
\href{https://github.com/amar-jay/amarcode/tree/main/crates/application/src}{\texttt{crates/application/src}}
& Sohbet ekranı, istem girişi, ajan seçimi, ayarlar, dosya ağacı ve diff
görüntüleme \\
Tauri katmanı &
\href{https://github.com/amar-jay/amarcode/tree/main/crates/application/src-tauri}{\texttt{crates/application/src-tauri}}
& Masaüstü penceresi, daemon bağlantısı, yerel dosya erişimi ve süreç
yönetimi \\
Amarcode daemon &
\href{https://github.com/amar-jay/amarcode/tree/main/crates/daemon}{\texttt{crates/daemon}}
& Kalıcı durum, RPC sunucusu, ACP oturumları, ajan süreçleri ve canlı
olaylar \\
Ortak protokol &
\href{https://github.com/amar-jay/amarcode/tree/main/crates/protocol}{\texttt{crates/protocol}}
& İstemci--daemon sözleşmesi, RPC türleri, olaylar ve TypeScript bağ
üretimi \\
ACP adaptörü &
\href{https://github.com/amar-jay/amarcode/tree/main/crates/amarcode-acp}{\texttt{crates/amarcode-acp}}
& OpenAI uyumlu sağlayıcıları ACP standardına bağlama ve çalışma alanı
araçları \\
Daemon registry &
\href{https://github.com/amar-jay/amarcode/tree/main/crates/daemon-registry}{\texttt{crates/daemon-registry}}
& Platforma uygun daemon sürümlerinin ve indirme bilgilerinin
yayımlanması \\
SQLite migration'ları &
\href{https://github.com/amar-jay/amarcode/tree/main/crates/daemon/migrations}{\texttt{crates/daemon/migrations}}
& Kalıcı veri şemasının sürümlü olarak oluşturulması ve güncellenmesi \\
Rapor görselleri &
\href{https://github.com/amar-jay/amarcode/tree/main/assets/diagrams}{\texttt{assets/diagrams}}
& Düzenlenebilir Mermaid kaynakları ile SVG ve PNG çıktıları \\
\end{longtable}

Rust çalışma alanının ana bileşenleri kök
\href{https://github.com/amar-jay/amarcode/blob/main/Cargo.toml}{\texttt{Cargo.toml}}
dosyasında, JavaScript çalışma alanı ve ortak komutlar ise kök
\href{https://github.com/amar-jay/amarcode/blob/main/package.json}{\texttt{package.json}}
dosyasında tanımlanmıştır.

\subsection{Ek B --- Genel Sistem
Mimarisi}\label{ek-b--genel-sistem-mimarisi}

\pandocbounded{\includegraphics[keepaspectratio,alt={Amarcode genel sistem mimarisi}]{/home/manan/Desktop/code/amarcode/assets/diagrams/system-architecture.png}}

\emph{Ek Şekil B.1 --- Kullanıcı, React/Tauri masaüstü uygulaması,
bağımsız daemon, SQLite, ACP ajanı, ajan kayıt sistemi ve çalışma alanı
arasındaki ilişkiler.}

Mimaride masaüstü arayüzü kalıcı verinin sahibi değildir. Kullanıcı
komutları Tauri köprüsü üzerinden yerel TCP JSON-line RPC bağlantısına
aktarılır. Daemon'ın servis katmanı iş akışlarını yönetir, kalıcı
kayıtları SQLite'a yazar ve yalnızca başarılı kayıttan sonra istemcilere
olay gönderir. ACP istemcisi kodlama ajanını ayrı bir süreç olarak
başlatır ve standart giriş/çıkış üzerinden JSON-RPC iletişimi kurar. Bu
ayrım, masaüstü penceresi kapansa bile görevlerin arka planda
sürdürülebilmesini sağlar.

Diyagramın düzenlenebilir kaynağı:
\href{https://github.com/amar-jay/amarcode/blob/main/assets/diagrams/system-architecture.mmd}{\texttt{system-architecture.mmd}}

\subsection{Ek C --- Temel Çalışma
Akışları}\label{ek-c--temel-uxe7alux131ux15fma-akux131ux15flarux131}

\subsubsection{Ek C.1 --- Kullanıcı isteminin yaşam
döngüsü}\label{ek-c1--kullanux131cux131-isteminin-yaux15fam-duxf6nguxfcsuxfc}

\pandocbounded{\includegraphics[keepaspectratio,alt={Kullanıcı istemi yaşam döngüsü}]{/home/manan/Desktop/code/amarcode/assets/diagrams/prompt-lifecycle-sequence.png}}

\emph{Ek Şekil C.1 --- Kullanıcı isteminin React arayüzünden daemon'a ve
ACP ajanına ulaşması ile akışlı cevabın arayüze geri dönmesi.}

Bu akışta kullanıcı mesajı arayüzde iyimser olarak gösterilir. Daemon
önce çalışma ve kullanıcı mesajı kayıtlarını oluşturur, ardından ACP
oturumunu başlatır veya sürdürür. Ajan tarafından üretilen her mesaj
parçası kalıcı depoya yazıldıktan sonra \texttt{messagePartAdded}
olayıyla arayüze bildirilir. Dönüş tamamlandığında mesaj ve çalışma
durumu güncellenerek \texttt{turnUpdated} olayı yayımlanır.

\subsubsection{Ek C.2 --- Önce kaydet, sonra bildir
ilkesi}\label{ek-c2--uxf6nce-kaydet-sonra-bildir-ilkesi}

\pandocbounded{\includegraphics[keepaspectratio,alt={Önce kaydet olay akışı}]{/home/manan/Desktop/code/amarcode/assets/diagrams/store-first-event-flow.png}}

\emph{Ek Şekil C.2 --- Bir ACP sinyalinin sahiplik kontrolünden, SQLite
işleminden ve olay yayınından geçişi.}

Bu ilke, arayüzün veritabanında bulunmayan bir durumu görmesini
engeller. Kalıcı yazım başarısız olursa işlem geri alınır ve başarılı
bir ürün olayı yayımlanmaz. Uygulama yenilendiğinde veya bağlantı tekrar
kurulduğunda arayüz aynı kalıcı kayıtlardan yeniden oluşturulabilir.

\subsubsection{Ek C.3 --- ACP araç izni ve terminal
yönetimi}\label{ek-c3--acp-arauxe7-izni-ve-terminal-yuxf6netimi}

\pandocbounded{\includegraphics[keepaspectratio,alt={ACP araç izni ve terminal sıralaması}]{/home/manan/Desktop/code/amarcode/assets/diagrams/acp-tool-permission-sequence.png}}

\emph{Ek Şekil C.3 --- Modelin komut çağrısı, kullanıcı izni, terminal
oluşturma, çıktı alma, iptal ve kaynak temizleme sırası.}

Komutlar örtük bir kabuk üzerinden yorumlanmaz; çalıştırılabilir dosya
ve argümanlar ayrı tutulur. Çalışma dizini aktif proje klasörüyle
sınırlandırılır. Kullanıcı tarafından verilen izin tam komut, argüman
listesi ve çalışma dizini için geçerlidir. İşlem başarıyla tamamlansa,
hata verse veya iptal edilse bile terminal kaynağı serbest bırakılır.

\subsection{Ek D --- Daemon ve Registry Yaşam
Döngüleri}\label{ek-d--daemon-ve-registry-yaux15fam-duxf6nguxfcleri}

\subsubsection{Ek D.1 --- Daemon
durumları}\label{ek-d1--daemon-durumlarux131}

\pandocbounded{\includegraphics[keepaspectratio,alt={Daemon yaşam döngüsü}]{/home/manan/Desktop/code/amarcode/assets/diagrams/daemon-lifecycle-state.png}}

\emph{Ek Şekil D.1 --- Daemon kontrolü, indirme, doğrulama, kurulum,
başlatma, güncelleme, geri alma ve kaldırma durumları.}

\subsubsection{Ek D.2 --- Güncelleme ve geri
alma}\label{ek-d2--guxfcncelleme-ve-geri-alma}

\pandocbounded{\includegraphics[keepaspectratio,alt={Daemon güncelleme sıralaması}]{/home/manan/Desktop/code/amarcode/assets/diagrams/daemon-update-rollback-sequence.png}}

\emph{Ek Şekil D.2 --- Yeni daemon sürümünün geçici dosyaya indirilmesi,
doğrulanması, yerleştirilmesi ve sağlık kontrolü başarısızsa önceki
sürümün geri yüklenmesi.}

\subsubsection{Ek D.3 --- Ajan registry ve kurulum
akışı}\label{ek-d3--ajan-registry-ve-kurulum-akux131ux15fux131}

\pandocbounded{\includegraphics[keepaspectratio,alt={Ajan registry ve kurulum akışı}]{/home/manan/Desktop/code/amarcode/assets/diagrams/agent-registry-install-flow.png}}

\emph{Ek Şekil D.3 --- Ajan manifestlerinin eşitlenmesi, dağıtım
türlerinin işlenmesi, çalışma zamanı denetimi ve kurulum süreci.}

Registry eşitlemesi en iyi çaba yaklaşımıyla uygulanır. Ağ bağlantısı
yoksa daemon son başarılı checkout üzerinden çalışmaya devam eder.
Manifestler \texttt{bunx}, \texttt{uvx} veya platforma özgü ikili
dağıtımlara dönüştürülür. Bir ajanın katalogda bulunması ile yerel
makinede kullanılabilir olması ayrı durumlar olarak değerlendirilir.

\subsection{Ek E --- Oturum Yapılandırmasının
Eşitlenmesi}\label{ek-e--oturum-yapux131landux131rmasux131nux131n-eux15fitlenmesi}

\pandocbounded{\includegraphics[keepaspectratio,alt={Oturum yapılandırma eşitleme akışı}]{/home/manan/Desktop/code/amarcode/assets/diagrams/session-config-sync.png}}

\emph{Ek Şekil E.1 --- ACP ajanının bildirdiği yapılandırma seçenekleri
ile daemon, SQLite ve React/Jotai durumu arasındaki çift yönlü veri
akışı.}

Ajanın bildirdiği seçim ve boolean seçenekleri genel arayüz
kontrollerine dönüştürülür. İlk istemden önce seçilen değerler
bekletilerek \texttt{prompt} isteğiyle gönderilir. Aktif oturum
sırasında değiştirilen değerler \texttt{session/set\_config\_option}
üzerinden ajana aktarılır. Yapılandırmanın tam anlık görüntüsü SQLite'ta
saklanır; boş liste de önceki değerleri temizleyen geçerli bir anlık
görüntü olarak kabul edilir. Ajan tarafından gönderilen güncellemeler
yalnızca aktif çalışma sahipliği doğrulandıktan sonra arayüze uygulanır.

\subsection{Ek F --- Seçilmiş Uygulama
Ekranları}\label{ek-f--seuxe7ilmiux15f-uygulama-ekranlarux131}

Bu bölümdeki ekran görüntüleri uygulamanın 7 Eylül 2026 tarihindeki son
durumunu göstermektedir. Görseller, ilgili özelliklerin rapor boyunca
açıklanan geliştirme çalışmalarının sonucunu göstermesi amacıyla
eklenmiştir; ekran görüntüsünün çekildiği tarih, özelliğin
geliştirildiği staj günü olarak yorumlanmamalıdır.

\subsubsection{Ek F.1 --- Açık ve koyu
tema}\label{ek-f1--auxe7ux131k-ve-koyu-tema}

\pandocbounded{\includegraphics[keepaspectratio,alt={Amarcode açık ve koyu tema}]{/home/manan/Desktop/code/amarcode/assets/amarcode-theme-split.png}}

\emph{Ek Şekil F.1 --- Amarcode'un açık ve koyu temalarının
karşılaştırmalı görünümü.}

\subsubsection{Ek F.2 --- Yeni sohbet
ekranı}\label{ek-f2--yeni-sohbet-ekranux131}

\pandocbounded{\includegraphics[keepaspectratio,alt={Amarcode yeni sohbet ekranı}]{/home/manan/Desktop/code/amarcode/assets/screenshots/Screenshot From 2026-09-07 13-20-36.png}}

\emph{Ek Şekil F.2 --- Çalışma alanı, ajan ve yapılandırma seçiminin
ardından yeni görev başlatılan sade başlangıç ekranı.}

\subsubsection{Ek F.3 --- Dosya ağacı ve kaynak
görüntüleyici}\label{ek-f3--dosya-aux11facux131-ve-kaynak-guxf6ruxfcntuxfcleyici}

\pandocbounded{\includegraphics[keepaspectratio,alt={Dosya ağacı ve kaynak görüntüleyici}]{/home/manan/Desktop/code/amarcode/assets/screenshots/Screenshot From 2026-09-07 13-18-38.png}}

\emph{Ek Şekil F.3 --- Aktif sohbetin yanında açılan çalışma alanı dosya
ağacı ve kaynak kodu görünümü.}

\subsubsection{Ek F.4 --- Diff inceleme}\label{ek-f4--diff-inceleme}

\pandocbounded{\includegraphics[keepaspectratio,alt={Amarcode diff inceleme ekranı}]{/home/manan/Desktop/code/amarcode/assets/screenshots/Screenshot From 2026-09-07 13-24-09.png}}

\emph{Ek Şekil F.4 --- Ajan tarafından oluşturulan değişikliklerin
eklenen ve silinen satırlar hâlinde incelenmesi.}

\subsubsection{Ek F.5 --- Ajan seçimi}\label{ek-f5--ajan-seuxe7imi}

\pandocbounded{\includegraphics[keepaspectratio,alt={Amarcode ajan seçimi}]{/home/manan/Desktop/code/amarcode/assets/screenshots/Screenshot From 2026-09-07 11-21-40.png}}

\emph{Ek Şekil F.5 --- Registry tarafından sağlanan ajanlar ile
kullanılabilirlik ve kurulum durumları.}

\subsubsection{Ek F.6 --- Model ve yapılandırma
seçenekleri}\label{ek-f6--model-ve-yapux131landux131rma-seuxe7enekleri}

\pandocbounded{\includegraphics[keepaspectratio,alt={Ajan yapılandırma seçenekleri}]{/home/manan/Desktop/code/amarcode/assets/screenshots/Screenshot From 2026-09-07 11-21-49.png}}

\emph{Ek Şekil F.6 --- Seçilen ACP ajanının bildirdiği model ve oturum
yapılandırma kontrolleri.}

\subsection{Ek G --- İstemci--Daemon RPC
Özeti}\label{ek-g--iux307stemcidaemon-rpc-uxf6zeti}

\begin{longtable}[]{@{}ll@{}}
\toprule\noalign{}
RPC yöntemi & Amaç \\
\midrule\noalign{}
\endhead
\bottomrule\noalign{}
\endlastfoot
\texttt{health} & Daemon durumu, sürümü, protokol sürümü ve dinleme
adresini döndürür \\
\texttt{version} & Daemon ve ortak protokol sürümünü bildirir \\
\texttt{subscribe\_events} & Sohbet, çalışma veya oturum filtresiyle
canlı olay akışına abone olur \\
\texttt{list\_agents} & Registry'den elde edilen ajanları ve yerel
kullanılabilirliklerini listeler \\
\texttt{install\_agent} & Manifestte belirtilen dağıtım yöntemini
kullanarak ajanı kurar \\
\texttt{authenticate\_agent} & Kimlik doğrulaması gerektiren ajan için
ACP doğrulama akışını başlatır \\
\texttt{create\_chat} & Çalışma alanına bağlı kalıcı sohbet oluşturur \\
\texttt{list\_chats} & İsteğe bağlı çalışma alanı filtresiyle sohbetleri
getirir \\
\texttt{get\_chat} & Sohbeti, mesajları ve mesaj parçalarını getirir \\
\texttt{delete\_chat} & Sohbet ve ilişkili kalıcı kayıtları temizler \\
\texttt{prompt} & Kullanıcı istemini, ekleri ve başlangıç
yapılandırmasını ajan çalışmasına aktarır \\
\texttt{set\_session\_config\_option} & Aktif ACP oturumundaki bir
yapılandırma değerini değiştirir \\
\texttt{cancel} & Sohbetteki aktif ajan dönüşünü iptal eder \\
\texttt{respond\_permission} & Ajanın bekleyen izin isteğine cevap
verir \\
\texttt{respond\_input} & Ajanın kullanıcıdan istediği ek bilgiye cevap
verir \\
\end{longtable}

RPC türlerinin güncel tanımları
\href{https://github.com/amar-jay/amarcode/blob/main/crates/protocol/src/rpc.rs}{\texttt{crates/protocol/src/rpc.rs}}
dosyasında bulunmaktadır.

\subsection{Ek H --- Canlı Olay
Türleri}\label{ek-h--canlux131-olay-tuxfcrleri}

\begin{longtable}[]{@{}ll@{}}
\toprule\noalign{}
Olay & Arayüzdeki anlamı \\
\midrule\noalign{}
\endhead
\bottomrule\noalign{}
\endlastfoot
\texttt{chatUpdated} & Sohbet başlığı veya sohbet listesi değişmiştir \\
\texttt{runUpdated} & Ajan çalışmasının başlangıç, çalışma, bitiş veya
hata durumu değişmiştir \\
\texttt{turnUpdated} & Belirli kullanıcı isteminin dönüş durumu
değişmiştir \\
\texttt{contextRestoration} & Önceki ACP oturumu veya kalıcı konuşma
geçmişi geri yüklenmektedir \\
\texttt{messageUpdated} & Bir mesajın akış veya tamamlanma durumu
değişmiştir \\
\texttt{messagePartAdded} & Metin, düşünme, araç veya diğer
yapılandırılmış mesaj parçası eklenmiştir \\
\texttt{approvalRequired} & Ajan bir işlem için kullanıcı izni
beklemektedir \\
\texttt{questionRequired} & Ajan görev için kullanıcıdan ek bilgi
istemektedir \\
\texttt{workspaceFilesChanged} & Çalışma alanındaki bir veya daha fazla
dosya değişmiştir \\
\texttt{agentConnectionChanged} & Ajan bağlantısı kurulmuş, kesilmiş
veya hata vermiştir \\
\texttt{agentAuthRequired} & Seçilen ajan kimlik doğrulaması
gerektirmektedir \\
\texttt{sessionConfigUpdated} & Oturuma ait yapılandırma seçenekleri
güncellenmiştir \\
\end{longtable}

Olay türlerinin güncel tanımları
\href{https://github.com/amar-jay/amarcode/blob/main/crates/protocol/src/events.rs}{\texttt{crates/protocol/src/events.rs}}
dosyasında bulunmaktadır.

\subsection{Ek I --- Test ve Doğrulama
Matrisi}\label{ek-i--test-ve-doux11frulama-matrisi}

\begin{longtable}[]{@{}lll@{}}
\toprule\noalign{}
Alan & Doğrulama yaklaşımı & İlgili kaynak \\
\midrule\noalign{}
\endhead
\bottomrule\noalign{}
\endlastfoot
Ön yüz canlı durumu & Mesaj parçaları, eski çalışma olayları ve sohbet
geçişleri için durum testleri &
\href{https://github.com/amar-jay/amarcode/blob/main/crates/application/src/state/daemon-events.test.ts}{\texttt{daemon-events.test.ts}} \\
Oturum yapılandırması & Tam ve boş anlık görüntü, ajan değişimi ve
bekleyen ilk istem değerleri &
\href{https://github.com/amar-jay/amarcode/blob/main/crates/application/src/state/session-config.test.ts}{\texttt{session-config.test.ts}} \\
Canlı sohbet & İyimser mesaj, akışlı içerik ve tamamlanma davranışları &
\href{https://github.com/amar-jay/amarcode/blob/main/crates/application/src/state/live-chat.test.ts}{\texttt{live-chat.test.ts}} \\
İzin modu & İzin düzeyleri ve kullanıcı tercihleri &
\href{https://github.com/amar-jay/amarcode/blob/main/crates/application/src/state/permission-mode.test.ts}{\texttt{permission-mode.test.ts}} \\
ACP protokolü & Oturum, mesaj kimliği, iptal, araç ve terminal mesaj
sıraları &
\href{https://github.com/amar-jay/amarcode/blob/main/crates/amarcode-acp/tests/protocol_transcript.rs}{\texttt{protocol\_transcript.rs}} \\
Daemon dikey dilimi & İstemden veritabanına, ACP'ye ve canlı olaya kadar
uçtan uca akış &
\href{https://github.com/amar-jay/amarcode/blob/main/crates/daemon/tests/vertical_slice.rs}{\texttt{vertical\_slice.rs}} \\
Ortak protokol & Üretilmiş TypeScript bağlarının Rust tanımlarıyla
güncel kalması &
\href{https://github.com/amar-jay/amarcode/blob/main/crates/protocol/src/bin/generate-types.rs}{\texttt{generate-types.rs}} \\
Registry Worker & Kayıt API davranışı ve sürüm verisi &
\href{https://github.com/amar-jay/amarcode/tree/main/crates/daemon-registry}{\texttt{daemon-registry}} \\
Masaüstü dağıtımı & GitHub Actions üzerinde platform derleme ve
paketleme &
\href{https://github.com/amar-jay/amarcode/tree/main/.github/workflows}{\texttt{.github/workflows}} \\
\end{longtable}

Projede tanımlanan temel doğrulama komutları şunlardır:

\begin{Shaded}
\begin{Highlighting}[]
\ExtensionTok{bun}\NormalTok{ run app:test}
\ExtensionTok{bun}\NormalTok{ run app:lint}
\ExtensionTok{bun}\NormalTok{ run tsc}
\ExtensionTok{bun}\NormalTok{ run worker:test}
\ExtensionTok{cargo}\NormalTok{ test }\AttributeTok{{-}{-}workspace}
\ExtensionTok{cargo}\NormalTok{ clippy }\AttributeTok{{-}{-}workspace} \AttributeTok{{-}{-}all{-}targets} \AttributeTok{{-}{-}all{-}features} \AttributeTok{{-}{-}} \AttributeTok{{-}D}\NormalTok{ warnings}
\ExtensionTok{cargo}\NormalTok{ fmt }\AttributeTok{{-}{-}all} \AttributeTok{{-}{-}check}
\end{Highlighting}
\end{Shaded}

\subsection{Ek J --- Git Tabanlı Geliştirme
Özeti}\label{ek-j--git-tabanlux131-geliux15ftirme-uxf6zeti}

İncelenen Git geçmişi 7 Ağustos 2026 ile 7 Eylül 2026 arasındaki
\textbf{114 commit} üzerinden değerlendirilmiştir. Commit tarihleri
günlük rapor başlıklarıyla bire bir eşleştirilmemiş, proje gelişiminin
teknik sırası korunarak 18 staj gününe dağıtılmıştır.

\begin{longtable}[]{@{}ll@{}}
\toprule\noalign{}
Geliştirme dönemi & Öne çıkan çalışmalar \\
\midrule\noalign{}
\endhead
\bottomrule\noalign{}
\endlastfoot
İlk prototip & İstem girişi, çalışma ekranı, TCP daemon ve test
istemcisi \\
Arayüz--daemon birleşimi & RPC, iyimser mesajlar, olay tabanlı yükleme
ve Jotai durum yönetimi \\
Güvenilirlik & Eş zamanlı istem kilitleri, eski olay koruması, abonelik
ve yeniden bağlantı \\
Ortak protokol & Rust türlerinin merkezîleştirilmesi, TypeScript üretimi
ve sürüm el sıkışması \\
Daemon dağıtımı & Kullanıcı servisi, tek örnek kilidi, indirme,
güncelleme, geri alma ve kaldırma \\
Çalışma alanı deneyimi & Dosya ağacı, arama, ekler, diff görüntüleyici
ve dosya bağlantıları \\
ACP uyumluluğu & Typed ACP, iptal, araç çağrıları, terminal ve izin
yönetimi \\
Registry sistemi & Ajan keşfi, manifest ayrıştırma, kullanılabilirlik ve
uçtan uca kurulum \\
Son entegrasyon & Dinamik oturum seçenekleri, kimlik doğrulama, testler
ve arayüz düzenlemeleri \\
\end{longtable}

Git geçmişi:
\href{https://github.com/amar-jay/amarcode/commits/main/}{Amarcode
commit listesi}

\subsection{Ek K --- Bilinen Sınırlamalar ve Gelecek
Çalışmalar}\label{ek-k--bilinen-sux131nux131rlamalar-ve-gelecek-uxe7alux131ux15fmalar}

\begin{itemize}
\tightlist
\item
  ACP adaptöründe oturum listeleme, silme ve tam kalıcı devam ettirme
  yetenekleri genişletilebilir.
\item
  Görsel, ses ve ek MCP özellikleri desteklenen ajan yeteneklerine göre
  eklenebilir.
\item
  Erişilebilirlik kontrolleri ve klavye ile tam gezinme için daha
  kapsamlı testler uygulanabilir.
\item
  Windows, Linux ve macOS paketleri için daha geniş otomatik uçtan uca
  test matrisi hazırlanabilir.
\item
  Daemon güncelleme ve bağlantı hataları için kullanıcıya sunulan
  tanılama bilgileri geliştirilebilir.
\item
  Çok büyük çalışma alanlarında dosya ağacı, arama ve diff görünümü için
  performans ölçümleri yapılabilir.
\item
  Wayland gibi platforma özgü pencere efektleri kontrollü özellik
  algılama ile ele alınabilir.
\item
  Registry manifestleri için imza doğrulama ve daha ayrıntılı güven
  zinciri değerlendirilebilir.
\end{itemize}

\subsection{Ek L --- İlgili Belgeler}\label{ek-l--iux307lgili-belgeler}

\begin{itemize}
\tightlist
\item
  \href{https://github.com/amar-jay/amarcode/tree/main/reports}{18
  günlük raporların bulunduğu klasör}
\item
  \href{https://github.com/amar-jay/amarcode/blob/main/reports/KAYNAKCA.md}{Toplu
  kaynakça}
\item
  \href{https://github.com/amar-jay/amarcode/blob/main/README.md}{Amarcode
  README}
\item
  \href{https://github.com/amar-jay/amarcode/blob/main/crates/daemon/README.md}{Daemon
  teknik belgesi}
\item
  \href{https://github.com/amar-jay/amarcode/blob/main/crates/amarcode-acp/README.md}{ACP
  adaptörü teknik belgesi}
\end{itemize}

\end{document}
